% ===================================================================
%  TABELO N-RO 13 — La hotelo. — La restoracio. — La kafejo.
%  Traduction 2026 d'apres texto/io, controlee sur texto/fr et
%  texto/es. Consigne : voir texto/eo/10-tabelo-01.tex.
%
%  Deux notes, marquees toutes deux « (*) », et deux appels : celui du
%  bloc 01-2 (la chambre) et celui du bloc 09-1 (le menu). L'ordre les
%  departage.
% ===================================================================

%%K t13-noto-1 noto
\VUnotes{143.2mm}{%
\textit{\VUgras{(*) Pri la detaloj de ĉambro, vidu la tabelon N\textsuperscript{o}. 6.}}%
}

%%K t13-apar-1 apar
\VUsaut{3.0mm}
\VUtitre{15.0pt}{60}{TABELO N\textsuperscript{o} 13}
\VUsaut{1.6mm}
\VUpk{10.2pt}{40}{La hotelo. --- La restoracio.}
\VUsaut{2.0mm}
\VUpk{10.2pt}{40}{La kafejo.}
\VUsaut{3.4mm}

%%K t13-01-1 p
1. --- Alveninte hieraŭ vespere en Rojanon, mi ekloĝis ĉe « \textit{Hotelo de
Oceano} » kiu estis rekomendita al mi de amiko.

%%K t13-01-2 p
Oni provizis al mi belan \VUgras{ĉambron} \textsuperscript{(2)}, sur la dua
\VUgras{etaĝo} \textsuperscript{(1)} kie mi tre bone dormis (*).

%%K t13-02-1 p
2. --- Ĉi-matene, elirante el mia ĉambro, en kiun la \VUgras{servistino}
\textsuperscript{(3)} estis alportinta la matenmanĝeton, mi eniris la
\VUgras{fumejon} \textsuperscript{(5)} kiu ankaŭ estas uzata kiel
\VUgras{legejo} \textsuperscript{(4)} por legi la \VUgras{ĵurnalojn}
\textsuperscript{(6)}, fumante \VUgras{cigaredon} \textsuperscript{(8)}. Post
legi, mi malsupreniris per la \VUgras{lifto} \textsuperscript{(9)} kaj iris
promeni tra la urbo.

%%K t13-03-1 p
3. --- Ĉar mi estis forgesinta demandi la \VUgras{oficiston}
\textsuperscript{(12)} de la hotelo pri kiu horo oni donas la manĝojn sur la
\VUgras{komuna tablo} \textsuperscript{(11)}, mi frue revenis kaj profitis tion
por iri bani min en la \VUgras{banejo} \textsuperscript{(10)} de la hotelo.
Poste mi iris denove al la \VUgras{kasejo} \textsuperscript{(15)} por telefoni
al unu el miaj amikoj. Sed la \VUgras{telefonilo} \textsuperscript{(13)} estis
okupata; vidante mian embarason la \VUgras{kasistino} \textsuperscript{(14)},
kiu estis registranta \VUgras{fakturon} \textsuperscript{(17)} sur la
\VUgras{vojaĝantregistro} \textsuperscript{(16)} petis la \VUgras{pordiston}
\textsuperscript{(18)} sendi la \VUgras{grumon} \textsuperscript{(19)} fari mian
komision \VUgras{bicikle} \textsuperscript{(20)}.

%%K t13-04-1 p
4. --- Mi dankis ŝin kaj eliris por doni gratifiketon al la \VUgras{portisto}
\textsuperscript{(24)} (la penlaboristo) kiu alportis el la stacidomo sur sia
\VUgras{manveturilo} \textsuperscript{(25)} mian \VUgras{ĉapelskatolon}
\textsuperscript{(26)}, \VUgras{mian valizon} \textsuperscript{(27)} kaj
\VUgras{mian vojaĝsakon} \textsuperscript{(28)}. La \VUgras{leteristo}
\textsuperscript{(21)} ĵus alveninta, donis al mi \VUgras{leteron}
\textsuperscript{(22)}.

%%K t13-05-1 p
5. --- Mi restis ankoraŭ kelkatempe sur la sojlo de la pordo, apud la
\VUgras{leterskatolo} \textsuperscript{(23)}, rigardonte la cirkuladon sur la
strato. La \VUgras{konduktoro} \textsuperscript{(30)} de la \VUgras{omnibuso}
\textsuperscript{(29)} estis ŝarĝanta sur sian veturilon la \VUgras{kofron}
\textsuperscript{(31)} de \VUgras{vojaĝanto} \textsuperscript{(32)} kiun la
\VUgras{hotelmastro} \textsuperscript{(33)} estis adiaŭanta. Juna
\VUgras{telegrafisto} \textsuperscript{(35)} senhaste alportis
\VUgras{telegramon} \textsuperscript{(34)} por kliento de la hotelo;
\VUgras{turisto} \textsuperscript{(36)} kun sia \VUgras{fotografaparato}
\textsuperscript{(38)} ŝultre, konsultis sian \VUgras{gvidlibron}
\textsuperscript{(37)}.

%%K t13-tit-1 sub
\VUsaut{4.0mm}
\VUpk{10.2pt}{40}{Tagmanĝa tempo.}
\VUsaut{3.0mm}

%%K t13-06-1 p
6. --- Antaŭ la krado, senzorga \VUgras{komisiisto} \textsuperscript{(39)}
dormetis inter sia \VUgras{portohoko} \textsuperscript{(40)} kaj sia
\VUgras{ciraĵkesto} \textsuperscript{(41)}, dum juna \VUgras{(tolaĵ)lavistino}
\textsuperscript{(42)} kiu portis \VUgras{korbon} \textsuperscript{(43)} da
tolaĵaro, ridis pro la solena mieno de la \VUgras{tabloficisto}
\textsuperscript{(44)} kiu staris apud la pordo.

%%K t13-07-1 p
7. --- Vidante la \VUgras{kuiriston} \textsuperscript{(45)}, kiu eliris el sia
\VUgras{kuirejo} \textsuperscript{(47)} situanta ĉe la \VUgras{subteretaĝo}
\textsuperscript{(48)} por sendi \VUgras{lavknabon} \textsuperscript{(46)} fari
komision, mi memoris ke la horo de la tagmanĝo estas baldaŭa kaj mi eniris la
restoraciejon, trairante la korton en kiu \VUgras{motoristo}
\textsuperscript{(50)} garis sian \VUgras{aŭtomobilon} \textsuperscript{(49)},
dum \VUgras{maŝinisto} \textsuperscript{(52)} riparis, laŭ la indikoj de la
\VUgras{biciklisto} \textsuperscript{(53)}, \VUgras{petroltriciklon}
\textsuperscript{(51)} kiu ĵus estis akcidentita.

%%K t13-08-1 p
8. --- Mi demandis junan \VUgras{grumon} \textsuperscript{(54)} kiu portis
\VUgras{bierglasojn} \textsuperscript{(55)}, ĉu la komuna tablo estas sub la
verando, kaj pro lia jesa respondo, mi okupis lokon ĉe unu el la tabloj kaj
igis doni al mi bonegan manĝon.

%%K t13-noto-2 noto
\VUnotes{143.2mm}{%
\textit{\VUgras{(*) Per helpo de la manĝolisto enkondukita en la vortaron, la
instruisto rajtos facile diversigi la suban menuon.}}%
}

%%K t13-tit-2 sub
\VUsaut{4.0mm}
\VUpk{10.2pt}{40}{Bonega tagmanĝo.}
\VUsaut{3.0mm}

%%K t13-09-1 p
9. --- La kelnero alportis al mi la menuon (*) de la tagmanĝo kaj la vinliston.
Mi elektis kaj mendis kiel \VUgras{akcesoraĵon salikokojn} kun \VUgras{butero},
kaj, kiel enirmanĝon, \VUgras{tripojn laŭ la modo de Caen}. Mi ne manĝis
\VUgras{fiŝon}, sed petis porcion da ŝafida \VUgras{kruro}, kaj, kiel
\VUgras{legomplado, spinacojn} kun kremo. Poste, ĉar neniu el la
\VUgras{intermanĝoj} plaĉis al mi, mi igis alporti diversajn desertojn,
\VUgras{fromaĝon}, \VUgras{fruktojn} kaj \VUgras{biskvitojn}. Mi krome aldonis
bonegan Saint-Émilion-\VUgras{vinon}, kaj tre kontenta pri mia manĝo mi
forlasis la tablon post donaci al la \VUgras{kelnero} \textsuperscript{(57)}
belan \VUgras{gratifiketon} \textsuperscript{(59)}.

%%K t13-10-1 p
10. --- Vidante min preta por foriri, la \VUgras{administranto}
\textsuperscript{(60)} proponis al mi iri al la balkono por admiri la belegan
panoramon de \VUgras{Oceano} \textsuperscript{(62)}. Malproksime oni ekvidis
fiŝ-\VUgras{barketojn} \textsuperscript{(63)}, \VUgras{velŝipojn}
\textsuperscript{(64)} kaj enorman \VUgras{paketŝipon} \textsuperscript{(65)}
kies nigra siluedo elstaras sur la blankaj \VUgras{nuboj}
\textsuperscript{(67)} de la \VUgras{horizonto} \textsuperscript{(66)}. Malpeza
venteto igis flirti la \VUgras{standardon} \textsuperscript{(70)} enigitan super
la \VUgras{ŝildo} \textsuperscript{(69)} de la \VUgras{halo}
\textsuperscript{(68)}.

%%K t13-tit-3 sub
\VUsaut{3.0mm}
\VUpk{10.2pt}{40}{Amuzoj.}
\VUsaut{3.0mm}

%%K t13-11-1 p
11. --- La spektaĵo estis tiel interesa, ke mi decidis supreniri morgaŭ sur la
terason de la kafejo kunportante \VUgras{binoklon} \textsuperscript{(71)} aŭ
\VUgras{lorneton} \textsuperscript{(72)}.

%%K t13-12-1 p
12. --- Forlasante la balkonon, mi iris al la kafejo de la hotelo por gluti
\VUgras{trinkaĵon} \textsuperscript{(73)}, ludante \VUgras{bilardpartion}
\textsuperscript{(75)}. Senhalte mi trairis la kafsalonon en kiu multaj
konsumantoj ludis per \VUgras{ŝako} \textsuperscript{(78)}, per \VUgras{damoj}
\textsuperscript{(79)}, per \VUgras{dominoŝtonoj} \textsuperscript{(80)}, per
\VUgras{triktrako} \textsuperscript{(81)} aŭ per \VUgras{kartoj}
\textsuperscript{(82)}, kaj mi rekte supreniris la \VUgras{bilardejon}
\textsuperscript{(74)}.

%%K t13-13-1 p
13. --- Kiam la partio estis finita, mi malsupreniris al la teretaĝo kaj petis
de la kelnero \VUgras{koverton} \textsuperscript{(84)}, \VUgras{paperon}
\textsuperscript{(83)}, \VUgras{poŝtmarkon} \textsuperscript{(85)},
\VUgras{plumon} \textsuperscript{(87)} kaj \VUgras{inkon}
\textsuperscript{(86)} por skribi leteron.

%%K t13-13-2 p
Poste mi vokis \VUgras{veturigiston} \textsuperscript{(92)} kies
\VUgras{veturilo} \textsuperscript{(88)} estis haltinta antaŭ la kafejo. Mi
demandis lin pri la prezo laŭ la horoj aŭ unu kuro, kaj, kiam ni konkordis pri
la prezo, li malfermis al mi la \VUgras{veturilpordon} \textsuperscript{(89)}
kaj instalis min. Li resupreniris sur la \VUgras{sidilon}
\textsuperscript{(91)}, rapide ekirigis la ĉevalojn per \VUgras{vipeto}
\textsuperscript{(93)} kaj ni foriris por promeno ĉarma.
\VUsaut{6.0mm}
\VUfilet{22mm}
